% Options for packages loaded elsewhere
% Options for packages loaded elsewhere
\PassOptionsToPackage{unicode}{hyperref}
\PassOptionsToPackage{hyphens}{url}
\PassOptionsToPackage{dvipsnames,svgnames,x11names}{xcolor}
%
\documentclass[
  english,
  11pt,
]{article}
\usepackage{xcolor}
\usepackage[margin=25mm]{geometry}
\usepackage{amsmath,amssymb}
\setcounter{secnumdepth}{-\maxdimen} % remove section numbering
\usepackage{iftex}
\ifPDFTeX
  \usepackage[T1]{fontenc}
  \usepackage[utf8]{inputenc}
  \usepackage{textcomp} % provide euro and other symbols
\else % if luatex or xetex
  \usepackage{unicode-math} % this also loads fontspec
  \defaultfontfeatures{Scale=MatchLowercase}
  \defaultfontfeatures[\rmfamily]{Ligatures=TeX,Scale=1}
\fi
\usepackage{lmodern}
\ifPDFTeX\else
  % xetex/luatex font selection
\fi
% Use upquote if available, for straight quotes in verbatim environments
\IfFileExists{upquote.sty}{\usepackage{upquote}}{}
\IfFileExists{microtype.sty}{% use microtype if available
  \usepackage[]{microtype}
  \UseMicrotypeSet[protrusion]{basicmath} % disable protrusion for tt fonts
}{}
\makeatletter
\@ifundefined{KOMAClassName}{% if non-KOMA class
  \IfFileExists{parskip.sty}{%
    \usepackage{parskip}
  }{% else
    \setlength{\parindent}{0pt}
    \setlength{\parskip}{6pt plus 2pt minus 1pt}}
}{% if KOMA class
  \KOMAoptions{parskip=half}}
\makeatother
% Make \paragraph and \subparagraph free-standing
\makeatletter
\ifx\paragraph\undefined\else
  \let\oldparagraph\paragraph
  \renewcommand{\paragraph}{
    \@ifstar
      \xxxParagraphStar
      \xxxParagraphNoStar
  }
  \newcommand{\xxxParagraphStar}[1]{\oldparagraph*{#1}\mbox{}}
  \newcommand{\xxxParagraphNoStar}[1]{\oldparagraph{#1}\mbox{}}
\fi
\ifx\subparagraph\undefined\else
  \let\oldsubparagraph\subparagraph
  \renewcommand{\subparagraph}{
    \@ifstar
      \xxxSubParagraphStar
      \xxxSubParagraphNoStar
  }
  \newcommand{\xxxSubParagraphStar}[1]{\oldsubparagraph*{#1}\mbox{}}
  \newcommand{\xxxSubParagraphNoStar}[1]{\oldsubparagraph{#1}\mbox{}}
\fi
\makeatother


\usepackage{longtable,booktabs,array}
\newcounter{none} % for unnumbered tables
\usepackage{calc} % for calculating minipage widths
% Correct order of tables after \paragraph or \subparagraph
\usepackage{etoolbox}
\makeatletter
\patchcmd\longtable{\par}{\if@noskipsec\mbox{}\fi\par}{}{}
\makeatother
% Allow footnotes in longtable head/foot
\IfFileExists{footnotehyper.sty}{\usepackage{footnotehyper}}{\usepackage{footnote}}
\makesavenoteenv{longtable}
\usepackage{graphicx}
\makeatletter
\newsavebox\pandoc@box
\newcommand*\pandocbounded[1]{% scales image to fit in text height/width
  \sbox\pandoc@box{#1}%
  \Gscale@div\@tempa{\textheight}{\dimexpr\ht\pandoc@box+\dp\pandoc@box\relax}%
  \Gscale@div\@tempb{\linewidth}{\wd\pandoc@box}%
  \ifdim\@tempb\p@<\@tempa\p@\let\@tempa\@tempb\fi% select the smaller of both
  \ifdim\@tempa\p@<\p@\scalebox{\@tempa}{\usebox\pandoc@box}%
  \else\usebox{\pandoc@box}%
  \fi%
}
% Set default figure placement to htbp
\def\fps@figure{htbp}
\makeatother



\ifLuaTeX
\usepackage[bidi=basic,shorthands=off]{babel}
\else
\usepackage[bidi=default,shorthands=off]{babel}
\fi
\ifLuaTeX
  \usepackage{selnolig} % disable illegal ligatures
\fi


\setlength{\emergencystretch}{3em} % prevent overfull lines

\providecommand{\tightlist}{%
  \setlength{\itemsep}{0pt}\setlength{\parskip}{0pt}}



 


\makeatletter
\@ifpackageloaded{caption}{}{\usepackage{caption}}
\AtBeginDocument{%
\ifdefined\contentsname
  \renewcommand*\contentsname{Table of contents}
\else
  \newcommand\contentsname{Table of contents}
\fi
\ifdefined\listfigurename
  \renewcommand*\listfigurename{List of Figures}
\else
  \newcommand\listfigurename{List of Figures}
\fi
\ifdefined\listtablename
  \renewcommand*\listtablename{List of Tables}
\else
  \newcommand\listtablename{List of Tables}
\fi
\ifdefined\figurename
  \renewcommand*\figurename{Figure}
\else
  \newcommand\figurename{Figure}
\fi
\ifdefined\tablename
  \renewcommand*\tablename{Table}
\else
  \newcommand\tablename{Table}
\fi
}
\@ifpackageloaded{float}{}{\usepackage{float}}
\floatstyle{ruled}
\@ifundefined{c@chapter}{\newfloat{codelisting}{h}{lop}}{\newfloat{codelisting}{h}{lop}[chapter]}
\floatname{codelisting}{Listing}
\newcommand*\listoflistings{\listof{codelisting}{List of Listings}}
\makeatother
\makeatletter
\makeatother
\makeatletter
\@ifpackageloaded{caption}{}{\usepackage{caption}}
\@ifpackageloaded{subcaption}{}{\usepackage{subcaption}}
\makeatother
\usepackage{bookmark}
\IfFileExists{xurl.sty}{\usepackage{xurl}}{} % add URL line breaks if available
\urlstyle{same}
% fallback for those not using the hyperref driver hyperxmp:
\makeatletter
\@ifundefined{xmpquote}{\newcommand{\xmpquote}[1]{#1}}{}
\makeatother
\hypersetup{
  pdftitle={Why checking a few points is not enough},
  pdflang={en},
  colorlinks=true,
  linkcolor={blue},
  filecolor={Maroon},
  citecolor={Blue},
  urlcolor={Blue},
  pdfcreator={LaTeX via pandoc}}


\title{Why checking a few points is not enough}
\usepackage{etoolbox}
\makeatletter
\providecommand{\subtitle}[1]{% add subtitle to \maketitle
  \apptocmd{\@title}{\par {\large #1 \par}}{}{}
}
\makeatother
\subtitle{An optional lesson about local success and global
requirements}
\author{}
\date{2026-09-12}
\begin{document}
\maketitle


\href{../paper-map/article.html}{Return to the paper map} · Optional
mathematical side article

A bridge can look sound where you inspect it and still have a weak
section between the inspection points. The same issue appears when we
calculate a profile at a finite set of positions.

Our earlier work on the vortex encountered exactly this kind of
difficulty. A very narrow parameter interval contained a failed
condition that a coarse scan could miss.

This article explains the lesson with a graph and a simple distinction:
\textbf{a set of successful checks is different from a bound that covers
a whole region}.

\subsection{A dip can hide between sample
points}\label{a-dip-can-hide-between-sample-points}

Imagine checking a road's height every hundred metres. All the measured
heights are positive. Does that prove the road never dips below a
particular level?

No.~A short dip could sit between two measurements. More samples help
only if they reach the troublesome scale or are combined with
information about how fast the road can change.

\begin{figure}

\centering{

\pandocbounded{\includegraphics[keepaspectratio,alt={A smooth curve contains a narrow negative dip near the middle; five widely spaced sample markers all remain positive.}]{figures/reader.png}}

}

\caption{\label{fig-sampling}All marked sample points lie above zero,
while a narrow dip between them falls below it. This invented graph
illustrates the difference between sampling and a guarantee over an
interval.}

\end{figure}%

The graph does not show the paper's vortex. It makes one logical point
visible: values at selected positions do not by themselves determine
what happens everywhere else.

\subsection{The vortex has conditions at its outer
edge}\label{the-vortex-has-conditions-at-its-outer-edge}

The inner profile must do more than satisfy its local equations. At its
outer edge, it needs the right shear and related inequalities so the
later construction can continue.

Here ``edge'' is the boundary of the inner calculation, not the boundary
of all the fluid. The joining annulus and exterior lie beyond it.

Our calculation therefore asked two different questions. Does the local
profile solve the leading equations accurately? And does it reach a
state suitable for the next construction step?

A ``yes'' to the first question does not force a ``yes'' to the second.

\subsection{What the narrow interval taught
us}\label{what-the-narrow-interval-taught-us}

For one of our finite parameter choices, a targeted scan found negative
angular shear in a narrow band. Increasing the accuracy and the number
of radial coefficients preserved the failure.

Choosing a larger scale parameter recovered the required signs at the
declared points. That supports those numerical observations. It does not
certify the condition between all the points.

The scale of the successful fine scan was striking: all 19 distinct
high-precision coordinates became the \textbf{same stored value} when
converted to ordinary double-precision numbers.

Asking a computer to evaluate 19 inputs is not really 19 different
checks if the computer has already rounded all the inputs to one value.
This is why the \href{../d7_conditioning/article.html}{article about
computer digits} matters.

The tiny width here is dimensionless. It is not a molecular size in a
physical fluid.

\subsection{A bound answers a different
question}\label{a-bound-answers-a-different-question}

Suppose we know the highest possible speed at which a road can rise or
fall over each stretch. Together with sampled heights, that information
may rule out a hidden deep dip.

More generally, a \textbf{bound} controls all allowed cases in a stated
region. It does not depend on happening to sample the worst point.

The study also constructed a conservative analytic bound for the axis
amplitude over a specified mathematical neighbourhood. That bound has a
broader scope than the sampled exit checks.

The source uses complex-variable estimates to control the profiles and
their derivatives. Those details are in the technical article. The
distinction to keep here is about coverage: a bound throughout a region
versus measurements at selected coordinates.
\href{https://cdn.openai.com/pdf/32d9f210-8b73-45e0-91bc-82a30aef8a9a/navier-stokes.pdf}{Paper,
Appendix B, especially (B.16)--(B.19)}

\subsection{Why very large numbers need a different
representation}\label{why-very-large-numbers-need-a-different-representation}

One amplitude depends on an exponential. Writing the enormous
normalization as an ordinary number is impractical.

We instead store its logarithm. If a quantity is \(10^{1000}\), storing
the exponent 1000 can be much more convenient than writing a one
followed by a thousand zeros. Logarithms also turn multiplication into
addition.

The actual calculation uses natural logarithms and extra arithmetic
precision. You do not need those details to see why changing the
representation helps.

\subsection{Where this fits in the bigger
picture}\label{where-this-fits-in-the-bigger-picture}

\textbf{This side study explains why reliable local calculations and
global construction conditions are separate tasks.} It supports the
inner-profile preparation in Appendix B, before the final joined field
is assembled.

It is optional on a first reading of the mechanism. Return to the
\href{../paper-map/article.html}{paper map}, or follow the main story
into \href{../convex-integration/article.html}{convex integration and
shear preparation}.

The \href{technical-article.html}{technical article},
\href{../../code/d6_global_axis/study.ipynb}{notebook}, and
\href{source-notes.md}{source notes} retain the actual parameter
choices, analytic bound, and numerical checks. The full joined/corrected
flow was not computed in this study.




\end{document}
